% ЕМТ 2023, XI клас — точен LaTeX препис от официалните файлове в Klasirane.com.
% Условия: https://klasirane.com/api/competitions/EMT/All/2023/11%20%D0%BA%D0%BB/probs
% SHA-256 (условия): 400dea00284d66fde519c47ba3bf142c680df4e0230efcadd155947ff5678da6
% Решения: https://klasirane.com/api/competitions/EMT/All/2023/11%20%D0%BA%D0%BB/sol
% SHA-256 (решения): 1f38d3ef7b13ee03f95d16c3fc5e45992fd000443a872639e431f11db03e30a2
% Точковите оценки, правописът, формулировките и видимите печатни особености са запазени от източника.
% В официалните файлове няма отделна чертежна фигура; не е добавяна такава.

Задача 11.1. Четворка $(a,b,c,d)$ от различни естествени числа се нарича $k$-хубава, ако са изпълнени следните две свойства:
\begin{itemize}
\item Измежду числата $a,b,c,d$ няма три, които да образуват (в някакъв ред) аритметична прогресия.
\item Измежду числата $a+b,a+c,a+d,b+c,b+d$ и $c+d$ има $k$, които образуват (в някакъв ред) аритметична прогресия.
\end{itemize}
а) Да се намери 4-хубава четворка.

б) Да се намери най-голямото $k$ за което съществува $k$-хубава четворка.

Решение. а) Четворката $(7,6,4,3)$ е хубава защото в нея няма три числа, които да образуват аритметична прогресия, а от шестте числа
$$
7+6=13,\quad 7+4=11,\quad 7+3=10,\quad 6+4=10,\quad 6+3=9,\quad 4+3=7
$$
числата 7, 9, 11, 13 образуват аритметична прогресия.

б) Без ограничение нека $a>b>c>d$. Тогава
$$
a+b>a+c>\max(a+d,b+c)>\min(a+d,b+c)>b+d>c+d.
$$
Да забележим, че ако:

1. $a+b,a+c$ и $a+d$ образуват аритметична прогресия, то $2(a+c)=(a+b)+(a+d)\Longleftrightarrow 2c=b+d$;

2. $a+b,a+c$ и $b+c$ образуват аритметична прогресия, то $2(a+c)=(a+b)+(b+c)\Longleftrightarrow 2b=a+c$;

3. $a+d,b+d$ и $c+d$ образуват аритметична прогресия, то $2(b+d)=(a+d)+(c+d)\Longleftrightarrow 2b=a+c$.

4. $b+c,b+d$ и $c+d$ образуват аритметична прогресия, то $2(b+d)=(b+c)+(c+d)\Longleftrightarrow 2c=b+d$.

И в четирите случая получаваме противоречие с условието на задачата.

От горното е ясно, че всичките 6 числа не могат да образуват аритметична прогресия.

Да допуснем, че 5 от тях образуват аритметична прогресия. Да забележим, че което и от числата $a+b,a+c,a+d,b+c,b+d$ и $c+d$ да изтрием, винаги се среща някоя от прогресийте 1., 2., 3., или 4.,противоречие.

От а) следва, че търсеното $k$ е 4.

Оценяване. (6 точки) а) за вярна 4-хубава четворка – 2 точки; б) за наредба на четирите числа и на получените шест сбора – 1 точка; за доказателство, че $k\leq4$ – 3 точки; частични резултати: за наблюдението, че някой от примерите 1., 2., 3. и 4. води до противоречие – 1 точка; за доказателство, че шестте числа не могат да образуват аритметична прогресия (еквивалентно на $k\leq5$) – 1 точка.

Задача 11.2. Върху страните $AB,BC$ и $AC$ на триъгълник $ABC$ са избрани съответно точки $C_1,A_1$ и $B_1$ така че $BA_1=BC_1$ и $CA_1=CB_1$. Правите $A_1C_1$ и $A_1B_1$ пресичат права през $A$, успоредна на $BC$, съответно в точки $P$ и $Q$. Ако описаните окръжности около триъгълниците $APC_1$ и $AQB_1$ се пресичат за втори път в точка $R$ върху отсечката $AA_1$, да се докаже, че точка $R$ лежи на вписаната в триъгълник $ABC$ окръжност.

Решение. От $PQ\parallel BC$ и $BA_1=BC_1$ получаваме
$$
\angle APC_1=\angle C_1A_1B=\angle A_1C_1B=\angle AC_1P.
$$
Тъй като четириъгълникът $APC_1R$ е вписан, имаме
$$
\angle A_1RC_1=\angle APC_1=\angle PC_1A=\angle PRA.
$$
Аналогично
$$
\angle CA_1B_1=\angle CB_1A_1=\angle AB_1Q=\angle AQB_1=\angle A_1RB_1=\angle ARQ.
$$
От $\angle B_1RC_1+\angle B_1A_1C_1=180^\circ$ следва, че $RC_1A_1B_1$ е вписан в окръжност $k$. Понеже $\angle C_1RA_1=\angle C_1A_1B=\angle A_1C_1B$ и $\angle B_1RA_1=\angle CA_1B_1=\angle CB_1A_1$, то $k$ се допира до страните на триъгълник $ABC$, т.е. $k$ е вписаната в $\triangle ABC$ окръжност.

Оценяване. (6 точки) За $\angle CA_1B_1=\angle CB_1A_1=\angle AB_1Q=\angle AQB_1=\angle A_1RB_1=\angle ARQ$ или съответното му – 2 точки; за $RC_1A_1B_1$ вписан – 2 точки; за извода, че $k$ е вписаната окръжност – 2 точки.

Задача 11.3. За естествено число $n$ са изпълнени следните сквойства:
\begin{itemize}
\item Числото $n+1$ се дели на 24.
\item Сборът от квадратите на всички делители на $n$ (включително 1 и самото $n$) се дели на 48.
\end{itemize}
Колко най-малко делители може да има $n$?

Решение. Тъй като $n+1$ се дели на 4, то $n$ не е точен квадрат. Следователно делителите на $n$ могат да бъдат разделени на двойки
$$
(a_0,b_0)=(1,n);(a_1,b_1),\ldots,(a_s,b_s),
$$
като броят на делителите на $n$ е $2(s+1)$. Тъй като 24 дели $n+1$, то всички делители на $n$ са нечетни и не се делят на 3. За всяко $i=0,1,\ldots,s$ имаме
$$
a_i+b_i=a_i+\frac{n}{a_i}=\frac{a_i^2-1+n+1}{a_i}
$$
и от $a_i$ нечетно, което не се дели на 3 следва, че $a_i^2-1$ се дели на 24. Следователно $a_i+b_i$ се дели на 24.

Сега от условието имаме, че
$$
\begin{aligned}
\sum_{i=0}^s(a_i^2+b_i^2)&=\sum_{i=0}^s((a_i+b_i)^2-2a_ib_i)\\
&=\sum_{i=0}^s(a_i+b_i)^2-2n(s+1)
\end{aligned}
$$
се дели на 48. Тъй като 48 дели $(a_i+b_i)^2$ и $n$ е нечетно, то 48 дели $2(s+1)$. Това означава, че $n$ има поне 48 делители.

Числото $n=23^{47}$ има исканите свойства, защото 24 дели $23^{47}+1$ и сборът от квадратите на делителите на $n$ е
$$
1+23^2+23^4+\cdots+23^{94}
$$
се дели на 48, защото $23^{2k}=529^k\equiv1\pmod{48}$.

Оценяване. (7 точки) За наблюдението, че всички делители на $n$ са нечетни и не се делят на 3 – 1 точка; за наблюдението, че $n$ не е точен квадрат и делителите му могат да се групират по двойки с произведение $n$ – 1 точка; за доказателство, че сборът на числата във всяка двойка се дели на 24 – 2 точки; за доказателство, че $n$ има поне 48 делители (т.е. 48 дели $2(s+1)$) – 1 точка; за намиране на число с 48 делители, което изпълнява условието – 2 точки.

Задача 11.4. Страната $A$ има $km$ града, а страната $B$ има $kn$ града $(k,m,n\in\mathbb N.)$. Всеки град от $A$ е свързан с двупосочна директна авиолиния с всеки град от $B$. Те се обслужват от $k$ авиокомпании (всяка авиолиния се обслужва само от една компания). Други авиолинии, освен посочените, няма. Докажете, че може да изберем авиокомпания и $m+n$ града, така че да е възможно да се придвижим между всеки два от избраните градове, ползвайки само авиолиниите на тази компания.

Решение. Лема 1. Нека $x,y$ са положителни реални числа, а $k,\ell,\ell\geq k$ са естествени числа. Реалните числа $x_i,y_i,i=1,2,\ldots,\ell$ удовлетворяват условията
$$
\begin{gathered}
x_i\geq0,\ y_i\geq0,\ x_i+y_i\leq\frac{x+y}{k},\quad i=1,2,\ldots,\ell;\\
\sum_{i=1}^{\ell}x_i=x,\qquad \sum_{i=1}^{\ell}y_i=y. \tag{1}
\end{gathered}
$$
Тогава е в сила неравенството
$$
\sum_{i=1}^{\ell}x_iy_i\leq\frac{xy}{k}.
$$
Равенството се достига само когато $x_i=x/k,y_i=y/k,i=1,2,\ldots,k;x_i=y_i=0,i>k$.

Доказателство. Да означим $f(x,y):=\sum_{i=1}^{\ell}x_iy_i$, където $x=(x_1,\ldots,x_\ell),y=(y_1,\ldots,y_\ell)$. Тъй като условията (1) определят компактно множество, функцията $f$ достига максималната си стойност върху него, да речем в точките $x_i',y_i'$. Можем да считаме, че $x_1'\geq x_2'\geq\cdots\geq x_\ell'$.

Ще докажем, че $y_i'$ са също в намаляваща последователност. Ако допуснем, че $y_i'<y_{i+1}'$, да разгледаме $x_i=x_{i+1}:=(x_i'+x_{i+1}')/2;y_i=y_{i+1}:=(y_i'+y_{i+1}')/2$. Тогава, (неравенство на Чебишев)
$$
x_iy_i+x_{i+1}y_{i+1}>x_i'y_i'+x_{i+1}'y_{i+1}'
$$
което противоречи на максималността на $x',y'$. По нататък, ако $x_1'+y_1'<(x+y)/k$ ние по аналогичен начин може да образуваме $x_1:=x_1'+\varepsilon,x_2:=x_2'-\varepsilon;y_1:=y_1'+\delta,x_2:=x_2'-\delta$ за подходящи $\varepsilon,\delta\geq0$ и да получим по-голяма стойност на $f$. Така че, $x_1'+y_1'=(x+y)/k$.

Нека $k'$ е най-голямото естествено число, за което $x_{k'}'>0$ и $y_{k'}'>0$. По същия начин, както по-горе, се вижда че $x_i'+y_i'=(x+y)/k,i=1,2,\ldots,k'$. Значи $k'\leq k$. Нека допуснем, че $k'<k$ и за определеност $y_i=0,i>k'$. Да модифицираме $x,y$ по следния начин. Полагаме $x_i:=x_i',y_i:=y_i',i=1,2,\ldots,k'-1;x_{k'}:=x_{k'}',y_{k'}:=y_{k'}'-\varepsilon,x_{k'+1}:=(x+y)/k,y_{k'+1}:=\varepsilon$.

За $i>k'+1$ числата $y_i$ са нули, а числата $x_i$ нямат значение, стига да се подчиняват на (1). Тъй като $x_k'<(x+y)/k$, лесно се вижда че $f(x,y)>f(x',y')$ което противоречи на максималността на $x',y'$.

И така, $k'=k$. Сега ще докажем, че $x_i'=x/k,y_i'=y/k,i=1,2,\ldots,k$. Да допуснем че това не е вярно и $j$ е първият индекс, за който $x_j'\ne x/k$, като нека за определеност нека $x_j'<x/k$. Тогава, ще съществува $i>j$ за което $x_i'>x/k$, което значи $x_i'>x_j'$ и значи последователността $x_1',x_2',\ldots,x_k'$ не е намаляваща, противоречие.

С това установихме, че $x_i'=x/k,y_i'=y/k,i=1,2,\ldots,k$. Тъй като $f(x',y')=kxy$, верността на Лема 1 е доказана.

Обратно към задачата. Броят всички авиолинии е $k^2mn$. Значи има авиокомпания, която обслужва поне $kmn$ авиолинии. Да премахнем всички останали авиолинии. Ще докажем, че в получения граф, нека бъде $K'$, има свързана компонента състояща се от поне $m+n$ върха.

Да допуснем противното. Нека свързаните компоненти на $K'$ са $G(A_i,B_i),i=1,2,\ldots,\ell$ и $|A_i|=m_i,|B_i|=n_i,i=1,2,\ldots,\ell$. Имаме
$$
\sum_{i=1}^{\ell}m_i=km,\qquad \sum_{i=1}^{\ell}n_i=kn,\qquad m_i+n_i<m+n
$$
Съгласно Лема 1,
$$
\sum_{i=1}^{\ell}m_in_i<kmn
$$
което противоречи на избора на авиолинията. И така, за поне едно $i$ е изпълнено $m_i+n_i\geq m+n$.

Оценяване. (7 точки) 2т. за стигане до неравенство от типа на Лема 1, 5т. за доказването му.
